% concatenated from Failure_of_weak_approximation_in_adjoint_groups.tex
\documentclass[11pt,letterpaper]{amsart}
\usepackage{amsmath}
\usepackage{amsfonts}
\usepackage{tikz-cd}
\usepackage{amsthm}
\usepackage{amssymb}
\usepackage{hyperref}
\usepackage{mathtools}
\usepackage{mathrsfs}
\usepackage{wasysym}
\usepackage{cite}

\theoremstyle{plain}
\newtheorem{theorem}{Theorem}           
\newtheorem{proposition}[theorem]{Proposition}
\newtheorem{corollary}[theorem]{Corollary}
\newtheorem{lemma}[theorem]{Lemma}
\newtheorem{claim}[theorem]{Claim}
\newtheorem{conjecture}[theorem]{Conjecture}

\theoremstyle{definition}
\newtheorem{definition}[theorem]{Definition}
\newtheorem{example}[theorem]{Example}
\newtheorem{question}[theorem]{Question}

\theoremstyle{remark}
\newtheorem*{remark}{Remark}
\newtheorem*{note}{Note}

\title{Failure of weak approximation in adjoint groups}
\author{Chayansudha Biswas}
\address{Department of Mathematics, University of Pennsylvania, Philadelphia, PA 19104-6395, USA}
\email{biswasch@sas.upenn.edu}
\subjclass{14G12, 11E57, 20G15, 14E08}
\keywords{weak approximation, adjoint groups, R-equivalence, rational varieties}

\begin{document}
    \maketitle

    \begin{abstract}
        Platonov in 1991 conjectured that adjoint groups are rational as varieties over arbitrary infinite fields, and as a consequence have weak approximation. The rationality part of the conjecture was disproved by Merkurjev in 1996, but the question about weak approximation remained open. We settle this in the negative.

    \end{abstract}

    \section*{Introduction}
    A generalization of the Chinese Remainder Theorem is the property of weak approximation of a field: if $S$ is any finite set of inequivalent discrete valuations on a field $F$, then the image of $F$ is dense in $\prod_{v\in S}F_v$ under the diagonal embedding, where $F_v$ denotes the completion of $F$ with respect to a valuation $v\in S$ \cite[Theorem 1.4]{MR1278263}. The topology considered here is the product adic topology. If $X$ is a variety defined over the field $F$, then one can ask if the image of $X(F)$ is dense in $\prod_{v\in S}X(F_v)$ under the diagonal embedding. If this is the case, then we say that $X$ \textit{has weak approximation with respect to} $S$. We say that $X$ has weak approximation if it has weak approximation with respect to every finite set of discrete valuations of $F$. All rational varieties have weak approximation, and hence failure of weak approximation is an obstruction to rationality.

    Weak approximation in the setting of algebraic groups is classical and well understood in many cases. For instance, it is well established that all semisimple simply connected and semisimple adjoint groups over global fields have weak approximation. The same statement is known to be false (due to work by Platonov in \cite{platonov1976reduced}) over arbitrary valued fields in the simply connected case, motivating a focus on adjoint groups instead. Platonov conjectured all semisimple adjoint groups over arbitrary fields are rational as varieties, and therefore have weak approximation \cite[\S 7.3]{MR1278263}. The rationality part of the conjecture was disproved by Merkurjev in \cite{merkurjev1996r}, where examples were given of adjoint groups that fail rationality. More examples of nonrational adjoint groups can be found in \cite{gille1997examples},  \cite{berhuy2004cohomological} and \cite{bhaskhar2014more}. 
    
    The obstruction to rationality in these examples is the failure of universal $R$-triviality: for a group $G$ defined over a field $F$, we say two points $x,y\in G(F)$ are \emph{$R$-equivalent} if there is a rational map $f\colon \mathbb{A}_F^1\dashrightarrow G$ defined at $0$ and $1$ such that $f(0) = x$ and $f(1) = y$. This is an equivalence relation on $G(F)$ and we denote by $G(F)/R$ the set of $R$-equivalence classes. We say the group $G$ is \emph{universally $R$-trivial} if $G(E)/R=1$ for all field extensions $E$ of $F$. This notion of $R$-equivalence was introduced by Manin in \cite{manin1986cubic}. Rationality of (the underlying variety of) a group implies universal $R$-triviality, and therefore Merkurjev constructed adjoint groups that fail universal $R$-triviality to obtain nonrational adjoint groups. 
    
    It was still not known, however, if there exist adjoint groups failing weak approximation. Over arbitrary fields, weak approximation has been shown for many cases of adjoint groups by Th\u{a}\'{n}g in \cite{thuang1996weak}, but the question remained open in general. To the author's knowledge, no further progress has been made towards the conjecture since.

    In the present article, we will settle the problem by constructing adjoint groups that fail weak approximation (see Theorem \ref{conjecture-disproved}, and see Example \ref{example} for an explicit example).
    For every global field $F$ of characteristic different from $2$ in which $-1$ is a square, we will give a recipe to construct infinitely many adjoint groups, defined over a suitable overfield of $F$, that fail weak approximation.

    To construct our counterexamples, we fix such a field $F$, and using a criterion of Merkurjev we first construct an adjoint group $G$ of type $D_n$ over~ $F$ such that $G(F)/R=1$ but $G(E)/R\neq 1$ for some extension $E$ over~ $F$. The main trick is to slightly modify the situation above to obtain a field $K$ over $F$ which admits a discrete valuation $v$ such that $G(K)/R=1$ and $G(K_v)/R\neq 1$. 
    We then use a result of Raghunathan to show that every~ $K_v$-rational point $g$ of $G$ has a $v$-adic neighborhood $U_g$ contained in the $R$-equivalence class of that point. In particular, we consider $g\in G(K_v)$ not in the $R$-equivalence class of $1$. Weak approximation in $G$ would imply that $U_g$ has nonempty intersection with $G(K)$. But $G(K)/R=1$ implies that points in this intersection are $R$-equivalent to the identity, which gives us a contradiction.

    \section{Constructing adjoint groups failing universal $R$-triviality}

    The aim of this section is the construction of adjoint groups that fail universal $R$-triviality. This is done in Proposition \ref{merk-psim}, using a criterion due to Merkurjev. Proposition \ref{merk-psim} will later be used in \S 2 to prove the main theorem. The following proposition is the main technical ingredient:
    
    \begin{proposition}\label{main-prop-sec1}
    Let $F$ be any global field with $\operatorname{char}F\neq 2$ in which $-1$ is a square. Then there is a quaternion algebra $Q$ over $F$ and an element~ $d\in~ F^{\times}\setminus F^{\times 2}$ such that $d\in -\operatorname{Nrd}(Q)$ and $d\notin\operatorname{Nrd}(Q^+)$. Here $Q^+$ denotes the subspace of $Q$ consisting of the pure quaternions, and $\operatorname{Nrd}\colon Q\to F$ is the reduced norm.
\end{proposition}

    For $F$ and $Q$ as in the proposition, let $A = M_n(Q)$, for $n\geq 3$ an odd integer. Then the hypothesis of \cite[Proposition 11]{merkurjev1996r} is satisfied, and thus there exists an involution $\sigma$ on $A$ with $\operatorname{disc}(\sigma) = d\cdot F^{\times 2}$. In the proof of this proposition it is shown that the hypothesis of \cite[Theorem 2]{merkurjev1996r} is satisfied, and hence the $F$-group $G = \mathbf{PSim}_+(A,\sigma)$ is an adjoint group of type $D_n$ failing universal $R$-triviality. (See \cite[p.2]{berhuy2004cohomological} for a definition of this group.)

    Failure of universal $R$-triviality means that there exists a field $E$ over $F$ such that $G(E)/R\neq 1$. The field $E$ is constructed from an iterated finite tower of fields over $F$ that are either function fields of quadrics or function fields of Severi-Brauer varieties (see \cite[Theorem 2]{merkurjev1996r}). This ensures that $E$ is a separably generated extension over $F$ of finite transcendence degree. We will use this fact to prove in the next section that $G$ fails weak approximation over a suitable overfield of $F$. We summarize this in the following proposition:
    \begin{proposition}\label{merk-psim}
        For any global field $F$ with $\operatorname{char}F\neq 2$ in which $-1$ is a square, there is an adjoint $F$-group $G$ of classical type $D_n$, where $n$ is an odd integer $\geq 3$, and a separably generated extension $E$ over $F$ of finite transcendence degree such that $G(E)/R\neq 1$.
    \end{proposition}

    \begin{remark}
        One checks, by following Merkurjev's construction, that the transcendence degree of $E$ over $F$ can be bounded above by $3n+4$ in the above notation.
    \end{remark}

    The remainder of this section is devoted to the proof of Proposition \ref{main-prop-sec1}.

    \begin{lemma}\label{d-exists}
    Let $F$ be any global field of characteristic different from $2$, and denote by $\Omega_F$ the set of all places of $F$. Let $\Sigma$ be any finite subset of $\Omega_F$. Then there is an element $d\in F^{\times}\setminus F^{\times 2}$ such that $d\in F_v^{\times 2}$ for all $v\in \Sigma$.
\end{lemma}
\begin{proof}
    Let $w$ be a place of $F$ not lying in the set $\Sigma$, and consider any element $a = (a_v)_{v\in\Sigma\cup\{w\}}\in \prod_{v\in\Sigma\cup\{w\}}F_v$ satisfying: \begin{enumerate}
        \item $a_v = 1$ for $v\in \Sigma$;
        \item $a_w\in F_w^\times\setminus F_w^{\times 2}$.
    \end{enumerate}
    Since $F$ is a global field, $F_v^{\times 2}$ is a subgroup of finite index in $F_v$ for $v\in\Omega_F$. This, coupled with the fact that $\operatorname{char}F_v\neq 2$, implies that $F_v^{\times 2}$ is $v$-adically open (and closed) in $F_v^{\times}$. Therefore, there is a neighborhood $U\subset \prod_{v\in\Sigma\cup\{w\}}F_v$ containing $a$ such that for any $b = (b_v)\in U$, $b_v$ is a square in $F_v$ for $v\in \Sigma$ and $b_w$ is not a square in $F_w$.

    By weak approximation for $F$, there is an element $d\in F\cap U$. Therefore, $d$ is a square in $F_v$ for all $v\in \Sigma$, and $d$ is not a square in $F_w$. In particular, we conclude $d\notin F^{\times 2}$. This element $d$ satisfies the conditions in the lemma.
\end{proof}

Fix a global field $F$ with $\operatorname{char}F\neq 2$ in which $-1$ is a square.
Choose a finite subset $\Sigma$ of $\Omega_F$ of even cardinality containing some valuation $v$ such that the corresponding residue field $\kappa_v$ has characteristic different from $2$. It is always possible to do so, since for our choice of $F$, there are at most finitely many valuations having residue characteristic $2$.
Define $Q$ to be the quaternion algebra over $F$ that is ramified precisely at the places in $\Sigma$; by \cite[Theorem 2.7.5]{maclachlan2013arithmetic}, $Q$ is uniquely determined up to isomorphism. Fix also $d\in F^\times$ as in Lemma \ref{d-exists}. We will retain this notation for the remainder of this section.

\begin{lemma}\label{neg-norm}
    One has $d\in -\operatorname{Nrd}(Q)$.
\end{lemma}
\begin{proof}
    We will use the Hasse-Schilling-Maass Theorem \cite[Theorem 33.15]{reiner2003maximal}, which says that $x\in\operatorname{Nrd}(Q)$ if and only if $x\in\operatorname{Nrd}(Q_v)$ for all places $v$ of $F$. It therefore suffices to show that $-d\in\operatorname{Nrd}(Q_v)$ for each $v\in\Omega_F$.

   First suppose $v\notin \Sigma$, i.e., $Q$ is not ramified at $v$. Then, $Q_v\simeq \operatorname{M}_2(F_v)$, and $\operatorname{Nrd}(Q_v)\simeq F_v^\times$. Therefore, $-d\in\operatorname{Nrd}(Q_v)$.

   Next, assume $v\in \Sigma$. By construction, $d\in F_v^{\times 2}$. Let $c\in F_v$ be such that $c^2 = d$. Then, $-d = (ic)^2 = \operatorname{Nrd}(ic)\in\operatorname{Nrd}(Q_v)$, where $i\in F$ is such that $i^2 = -1$.
\end{proof}

\begin{lemma}\label{ab-vals}
    Given $v\in \Sigma$, there exist $a,b\in F$ with $Q\simeq (a,b)_F$ such that $v(b)=1$ and $v(a)=0$. 
\end{lemma}
\begin{proof}
    Let $(a,b)_F$ be any representation of $Q_F$. Let $\pi$ be a uniformizer corresponding to the place $v$. We may write $a = a_0\pi^{v(a)}, b= b_0\pi^{v(b)}$ for $v$-adic units $a_0,b_0$.
    We will first show that at least one of $v(a)$ and $v(b)$ is nonzero.
    
     Recall that for a place $u\in\Omega_F$ the Hilbert symbol at $u$ is defined as the function $h_u\colon F_u^\times\times F_u^\times\to \{\pm 1\}$, which sends the pair $(a,b)$ to $1$ if and only if the quaternion algebra $(a,b)_{F_u}$ is split, i.e., isomorphic to $\operatorname{M}_2(F_u)$.
     By \cite[12.4.8]{voight2021quaternion}, we have $h_v(a,b) = (-1)^{v(a)v(b)(q-1)/2}\left(\frac{a_0}{v}\right)^{v(b)}\left(\frac{b_0}{v}\right)^{v(a)}$, where $q$ is the cardinality of the residue field $\kappa_v$, and $\left(\frac{a_0}{v}\right)$ and $\left(\frac{b_0}{v}\right)$ denote the Legendre symbols. If $v(a) = 0 = v(b)$, then $h(a,b) = 1$, which would mean that $(a,b)_{F_v}$ is split, a contradiction to the fact that $Q$ is ramified at $v$.

    Since $(a,b)_F\simeq (b,a)_F$, we may assume without loss of generality that $v(b)\neq 0$. Throughout the rest of the proof, we will frequently use the fact that scaling $a$ or $b$ by the square of a unit in $F$ does not change the isomorphism class of $(a,b)_F$.
    If $v(b)<0$, then $v(\pi^{-2v(b)}b) = -v(b)>0$. Therefore, by replacing $b$ by $\pi^{-2v(b)}b$ if necessary, we have $Q\simeq (a,b)_F$ with $v(b)>0$.

    If $v(a)<0$, replace $a$ by $a\pi^{-2v(a)}$ to ensure $v(a)\geq 0$. Assume without loss of generality that $v(a)\leq v(b)$ (since $(a,b)_F\simeq (b,a)_F)$). Recall that we have already shown in the first part of the proof that the case $(a,b)_F\simeq (a_0,b_0)_F$ is not possible. So by scaling each argument by a suitable power of $\pi^{-2}$, we may assume that $(a,b)_F$ is isomorphic to either $(a_0,b_0\pi)_F$ or $(a_0\pi,b_0\pi)$.  In the first case, simply replace $a$ by $a_0$ and $b$ by $b_0\pi$. In the second case, use the isomorphism $(x,y)_F\simeq (-xy,y)_F$ to obtain $(a,b)_F\simeq (a_0\pi,b_0\pi)_F\simeq (-a_0b_0\pi^2,b_0\pi)_F\simeq (-a_0b_0,b_0\pi)_F$. Now replace $a$ by $-a_0b_0$ and $b$ by $b_0\pi$.
\end{proof}

\begin{lemma}\label{not-pure-quat-norm}
    The element $d\in F^\times$ is not the norm of a pure quaternion in $Q$.
\end{lemma}
\begin{proof}
    Fix a valuation $v\in \Sigma$ with uniformizer $\pi$ such that the residue field $\kappa_v$ has characteristic different from $2$. By Lemma \ref{ab-vals}, there exist $a,b\in F$ with $Q\simeq (a,b)_F$ and such that $v(a) = 0$ and $v(b)=1$.
    Suppose $d = \operatorname{Nrd}(xI + yJ + zK) = -ax^2-by^2+abz^2$ for some $x,y,z\in F$. By multiplying throughout by a suitable $n$-th power of $\pi^2$, we obtain an equation of the form
    \begin{equation}\label{normform}
        D = -aX^2-bY^2 + abZ^2,
    \end{equation} where $D = d\pi^{2n}$, and $X = x\pi^{n},Y = y\pi^{n}$ and $Z = z\pi^{n}$ are in $\mathcal{O}_v$ and at least one of them is not in $\mathfrak{m}_v$. 
    
    Reducing the equation modulo $\mathfrak{m}_v$, we obtain $D\equiv -aX^2(\operatorname{mod}\mathfrak{m}_v)$. First assume $X\in\mathcal{O}_v^\times$. Since $d$ is a square in $F_v$, so is $D$, and let $C\in F_v$ be such that $C^2 = D$. Then we find $(iC/X)^2\equiv a(\operatorname{mod}\mathfrak{m}_v)$, which would mean that $a$ is a square modulo $\mathfrak{m}_v$, which contradicts the fact that $Q_{F_v}$ is not split. (Here, as before, $i$ is a square root of $-1$ in $F$.)

    So $v(X)>0$, and either $Y$ or $Z$ has zero valuation. Suppose it is $Y$. The case $Z\in\mathcal{O}_v^\times$ is similar. Since $v(b) = 1$, let $b = b_0\pi$ for $b_0\in\mathcal{O}_v^\times$. The congruence $D\equiv -aX^2\operatorname{mod}\mathfrak{m}_v$ gives $D\equiv X^2\equiv 0\operatorname{mod}\mathfrak{m}_v$. Let $D = D_0\pi^k, X= X_0\pi^j$ for $D_0,X_0\in\mathcal{O}_v^\times$ and $k,j\geq 1$. On the other hand, since $D = d\pi^{2n}$, one has $k = v(D) = v(d) + 2n$, which is even because $d$ is a square in $F_v$ (and hence $v(d)\in 2\mathbb{Z}$). Therefore, $k\geq 2$. Substituting these in (\ref{normform}), we obtain:\begin{equation}
        D_0\pi^k = -aX_0^2\pi^{2j} - b_0\pi Y^2 + ab_0\pi Z^2.
    \end{equation}
    Dividing throughout by $\pi$, we get
    \begin{equation}
        D_0\pi^{k-1} = -aX_0^2\pi^{2j-1} - b_0Y^2 + ab_0Z^2;
    \end{equation}
    and reducing this modulo $\mathfrak{m}_v$ gives:\begin{equation}
        b_0Y^2\equiv ab_0Z^2\operatorname{mod}\mathfrak{m}_v.
    \end{equation}
    Since $b_0\in\mathcal{O}_v^\times$, we can multiply both sides by $b_0^{-1}$ to obtain $Y^2\equiv aZ^2\operatorname{mod}\mathfrak{m}_v$. Since $v(Y) = 0$ by assumption and since $v(a)=0$, $Z\not\equiv 0\operatorname{mod}\mathfrak{m}_v$ i.e. $Z\in\mathcal{O}_v^\times$ and so we find $(Y/Z)^2\equiv a\operatorname{mod}\mathfrak{m}_v$, which would mean that $a$ is a square modulo $\mathfrak{m}_v$, again a contradiction.
\end{proof}

Proposition \ref{main-prop-sec1} now follows from Lemmas \ref{d-exists}, \ref{neg-norm} and \ref{not-pure-quat-norm}.

\section{Failure of weak approximation}
Fix a global field $F$ having characteristic different from $2$ and in which $-1$ is a square. Fix $G$ and $E$ as in Proposition \ref{merk-psim}.

\begin{lemma}\label{efg}
        For $F,G$ and $E$ as above, one has $G(F)/R = 1$ and $G(E)/R\neq~ 1$.
    \end{lemma}
    \begin{proof}
         We only need to show $G(F)/R = 1$. Since $F$ is either a totally imaginary number field or of the form $\mathbb{F}_q(t)$ with $q\neq 2$, it has cohomological dimension $2$ (see \cite[Proposition 8.3.17, Theorem 10.1.11]{MR1737196}), and so we may use \cite[Theorem 3.7]{kulshrestha2008} to conclude $G(F)/R=~1$.
    \end{proof}

Our next goal is to show that the base change of $G$ to a suitable overfield of $F$ fails weak approximation. We begin with the following lemma:
    
    \begin{lemma}\label{val-fld}
        There is a field $K$ containing $F$, and a valuation $v$ on $K$ with corresponding completion $K_v$, such that $G(K)/R = 1$ and $G(K_v)/R\neq 1$.
    \end{lemma}
    \begin{proof}
        Choose a separating transcendence basis $\{e_1,\ldots,e_n\}$ of $E$ over $F$. Then the field $F' = F(e_1,\ldots,e_n)$ is a purely transcendental extension over $F$ while $E$ is a separable algebraic extension of $F'$.
        By the Primitive Element Theorem, $E = F'(\alpha)$ for some $\alpha\in E$.

        Let $f_\alpha\in F'[t]$ be the minimal polynomial of the element $\alpha$. Then we may endow the field $K = F'(t)$ with valuation $v$ induced by the prime ideal $(f_\alpha)$ in $F'[t]$. Observe that the corresponding residue field $\kappa_v$ is isomorphic to $E$. Denote by $K_v$ the completion of $K$ with respect to $v$. Observe that $G(K_v)/R\simeq G(\kappa_v)/R\simeq G(E)/R\neq 1$, where the first isomorphism comes from \cite[Theorem 8.14]{gille2025r}, applied with $A = \mathcal{O}_v\supseteq F$. On the other hand, $G(K)/R\simeq G(F')/R\simeq G(F)/R = 1$ by homotopy invariance (see \cite[p.1]{gille2010lectures}).
    \end{proof}

    The next two lemmas prepare the key step in our argument: we show that the $R$-equivalence class of any $g\in G(K_v)$ is $v$-adically open. This means, in particular, that a $K_v$-rational point $g$ lying outside the $R$-equivalence class of 1 cannot be approximated by points $R$-equivalent to $1$.

    \begin{lemma}\label{raghunathan}
        Let $k$ be a field and let $\mathcal{G}$ be a connected reductive group over $k$. Then for any $g\in \mathcal{G}(k)$ there is a $k$-rational variety $Y$ and a $k$-morphism $b\colon Y\to \mathcal{G}$ which is \'{e}tale over $g$.
    \end{lemma}
    \begin{proof}
        It suffices to prove the lemma for $g=1$. Indeed, once we know that $b\colon Y\to\mathcal{G}$ is \'{e}tale over $1$, we can compose it with left translation by any $g\in\mathcal{G}(k)$ to obtain a map $b_g\colon Y\to\mathcal{G}$ that is \'{e}tale over $g$.
        
        
        We begin with a result due to Raghunathan \cite[\S 1.2]{raghunathan1994principal} (see also \cite[Proposition 3.7(1)]{gille2026local}): Let $T$ be a maximal torus of $\mathcal{G}$. Then there is $n\in\mathbb{N}$ and a $k$-morphism of varieties $f:T^n = T\times \cdots\times T\to \mathcal{G}$ that is \'{e}tale over the identity in $\mathcal{G}$. This comes from the isomorphism $\operatorname{Lie}(\mathcal{G})(k) = \mbox{}^{g_1}\operatorname{Lie}(T)(k) + \mbox{}^{g_2}\operatorname{Lie}(T)(k) + \cdots + \mbox{}^{g_n}\operatorname{Lie}(T)(k)$, for elements $g_i\in \mathcal{G}(k)$.

        Now let $1\to S\to Q\to T\to 1$ be a flasque resolution of $T$, where $S$ is a flasque torus and $Q$ is a quasitrivial torus (see \cite[Th\'{e}or\`{e}me 2]{colliot1977r}). The torus $Q$ is a Zariski open subset of an affine space $\mathbb{A}^m$, where $m = \dim_k Q$. Denote by $W$ the Lie algebra of $Q$. Then $W(k)\simeq k^m$.
        The epimorphism $Q\to T$ from the flasque resolution above induces an epimorphism $h\colon Q^n\to T^n$ of $k$-varieties. The differential $dh\colon W(k)^n\simeq k^{mn}\to \operatorname{Lie}(T)(k)^n$ is a surjective homomorphism of vector spaces. Observe $\mathbb{A}^{mn}\simeq \operatorname{Spec(Sym(}(W(k)^n)^\lor))$, where $\operatorname{Sym}((W(k)^n)^\lor)$ is the symmetric algebra on the dual space of $W(k)^n$. Let $V$ be a subspace of $W(k)^n$ such that the restriction of $dh$ to $V$ is an isomorphism.
        Denote by $Z$ the affine variety $\operatorname{Spec(Sym}(V^\lor))$. Notice that $Z$ is in fact isomorphic to an affine space $\mathbb{A}^r\subset\mathbb{A}^{mn}$, where $r = \dim_k V$.
        Let $Y:= Q^n\cap Z$, where the intersection takes place in $\mathbb{A}^{mn}\simeq \operatorname{Spec(Sym(}(W(k)^n)^\lor))$. Then $Y$, being an open subset of the affine space $Z$ is rational as a variety, and $h|_Y\colon Y\to T^n$ is \'{e}tale over $1$.

        The composition of $h|_Y$ followed by $f$ is the required map $b\colon Y\to \mathcal{G}$.
    \end{proof}

    \begin{lemma}\label{r-eq-open}
        Let $\mathcal{G}$ be a connected reductive group over a complete valued field $k$. Then, for any $g\in \mathcal{G}(k)$, there is an adic neighborhood $U$ of $g$ which is contained in the $R$-equivalence class of $g$.
    \end{lemma}
    \begin{proof}
        By Lemma \ref{raghunathan}, there is a morphism $b\colon Y\to \mathcal{G}$ from a rational variety~ $Y$ which is \'{e}tale over $g\in G(k)$. By the inverse function theorem over complete valued fields (see \cite[Corollary 2.1.1(ii)]{igusa2000introduction}), $b$ is an adic homeomorphism (between sets of $K_v$-rational points) over an open neighborhood $U$ of $g$ in $\mathcal{G}(K_v)$.

        We claim that every $h\in U$ is $R$-equivalent to $g$. Indeed, let $x,y\in Y(k)$ be such that $b(x)=g$ and $b(y) = h$. Since $Y$ is a rational $k$-variety, it is $R$-trivial and hence there is a $k$-rational morphism $f\colon\mathbb{A}_k^1\dashrightarrow Y$ with $f(0) = x$ and $f(1)=y$. Then, the $k$-rational map $b\circ f\colon \mathbb{A}_k^1\to G$ satisfies $(b\circ f)(0)=g$ and $(b\circ f)(1)=h$.
    \end{proof}

    We are now ready to state and prove the main theorem:

    \begin{theorem}\label{conjecture-disproved}
    There exists an adjoint group $G$ defined over a  field $K$, with a discrete valuation $v$ and corresponding completion $K_v$, such that the embedding $G(K)\to G(K_v)$ does not have $v$-adically dense image. In particular, $G$ does not have weak approximation.
\end{theorem}
\begin{proof}
    Let $F$ be any global field of characteristic different from $2$ in which $-1$ is a square. Let $G$ be as in Proposition \ref{merk-psim} and let $K,v$ and $K_v$ be as in Lemma \ref{val-fld}.
    
    Let $g\in G(K_v)$ be a $K_v$-rational point that is not in the $R$-equivalence class of $1$. By Lemma \ref{r-eq-open}, there is a $v$-adically open neighborhood $U$ of $g$ in $G(K_v)$ that is contained in the $R$-equivalence class of $g$.

    If $G$ had weak approximation, then there would be an element $h\in G(K)\cap U\subset G(K_v)$. Hence, $h$ is $R$-equivalent to $g$. Since $G(K)/R=1$, the identity element $1$ would be in the $R$-equivalence class of the element $h$ in $G(K)$ and therefore also in $G(K_v)$. This is a contradiction since $g$ was chosen to be not $R$-equivalent to $1$.
\end{proof}

This answers the Platonov-Th\u{a}\'{n}g conjecture negatively. We conclude by providing an explicit counterexample:

\begin{example}\label{example}
    Consider $F = \mathbb{Q}(i)$ and the quaternion algebra $Q = (2,5)_F$ over $F$. This is the unique (up to isomorphism) quaternion algebra over $F$ ramified exactly at the places $(2+i)$ and $(2-i)$. One may check that $d=6$ satisfies the conditions in Proposition \ref{main-prop-sec1}. Then, by Proposition \ref{merk-psim}, the algebra $A = \operatorname{M}_3(Q)$ has an involution $\sigma$ such that $G = \textbf{PSim}_+(A,\sigma)$ is not universally $R$-trivial.

    Lemma \ref{val-fld} and Theorem \ref{conjecture-disproved} then imply that $G_K$ fails weak approximation for some $K$ over $F$. Following Merkurjev's construction in \cite[Theorem 2, Lemma 8]{merkurjev1996r}, one can find such a $K$ with  $\operatorname{tr.deg}_FK\leq 13$.

    More generally, for each $A_n = M_n(Q)$ with $n\geq 3$ odd, there is an involution $\sigma_n$ such that $G_n = \textbf{PSim}_+(A_n,\sigma_n)$ does not have weak approximation when base changed to some field $K$ of transcendence degree at most $3n+4$ over $F = \mathbb{Q}(i)$.
\end{example}


\section*{Acknowledgments}
I am indebted to my advisor, Julia Hartmann, for her guidance, careful reading of multiple drafts, and numerous suggestions that substantially improved this work. I also thank Danny Krashen for many insightful discussions and for recommending the use of Lemma \ref{raghunathan}. I am grateful to David Harbater for helpful conversations and to Igor Rapinchuk for thoughtful comments on the manuscript. I thank Mattie Ji and Xingyu Meng for assistance with preliminary computational work that helped guide some results. This material is based upon work supported in part by the U.S. National Science Foundation under Grant No. DMS-2102987.


    



     \bibliographystyle{alpha}
    \bibliography{bibliography}{}

\end{document}


